\documentclass{situationist-preprint}

\title{Errata to \enquote{On the Hard Problem of Consciousness}, Volume III}
\msid{MND-01-e}
\msfamily{MND-01}
\msdiscipline{Philosophy of Mind}
\mstype{Erratum}
\mskeywords{corrigendum; Russellian monism; protophenomenal properties; local supervenience}

\begin{document}
\maketitle

\begin{abstract}
We correct six errors in Volume~III of MND-01 (\emph{The Metaphysics of Russellian Monism}, Chapters 21--30): two misnumbered references, a lemma stated without a necessary side-condition, a terminological slip, and two citation errors. For each we give the printed text, the correction, and the results that depend on it. No theorem of the volume changes its statement except Lemma~24.3, whose corrected form still supports every result derived from it. The conclusion of the work is unaffected.
\end{abstract}

Volume~III develops the metaphysics of the Russellian monism whose truth, given materialism, Volume~I establishes (MND-01-a, Corollary~4.3). Pagination is continuous across the four volumes; page numbers below are those of Volume~III.

\section{Misnumbered references}

\paragraph{1. Chapter~22, proof of Theorem~22.5 (p.~803).} Theorem~22.5 states that the protophenomenal natures posited in Chapter~21 are physical in the sense of MND-02. The proof cites \enquote{Proposition~3.5 of MND-02-a} for the claim that a categorical realizer of a property playing a dynamical role may be a fundamental property of a foreseeable extension of current physics. The claim is clause~(b) of Definition~2.2 of MND-02-a; Proposition~3.5 there is the non-triviality result. The reference should read \enquote{Definition~2.2(b) of MND-02-a}. The second step of the proof, which uses clause~(c), is cited correctly. Theorem~22.5 stands as printed.

\paragraph{2. Chapter~26, statement of Corollary~26.4 (p.~981).} Corollary~26.4 is said to follow \enquote{from Theorem~4.2 of Volume~I and Theorem~22.5}. Theorem~4.2 of Volume~I is the disjunction that either materialism is false or Russellian monism is true; the result needed is its consequence, given materialism, that Russellian monism is true, which is Corollary~4.3 of Volume~I. The reference should read \enquote{Corollary~4.3 of Volume~I}. The proof, which uses materialism explicitly in its first line, is correct as printed.

\section{A missing side-condition}

\paragraph{3. Chapter~24, Lemma~24.3 (p.~887).} Chapter~24 argues that a subject's experience supervenes locally on the protophenomenal natures and structure of the arrangement that grounds it. Lemma~24.3 is printed as follows.

\begin{quote}\small
\textbf{Lemma 24.3.} Any two arrangements alike in the intrinsic natures of their fundamental constituents and in their structure are alike in the experiences they ground.
\end{quote}

As printed, the lemma is false under the maximality condition of Definition~24.1, on which an arrangement grounds a subject only if it is not a proper part of an arrangement that realizes the same integrated organization. An arrangement that grounds a subject when isolated may have an exact duplicate embedded in a larger integrated arrangement, and that duplicate grounds no subject of its own. Since maximality is an extrinsic condition, intrinsic likeness does not guarantee likeness in grounded experience. The lemma should read:

\begin{quote}\small
\textbf{Lemma 24.3 (corrected).} Any two arrangements alike in the intrinsic natures of their fundamental constituents and in their structure, each of which is maximal in the sense of Definition~24.1, are alike in the experiences they ground.
\end{quote}

The lemma is used twice. In the proof of Theorem~24.6 (local supervenience of a normal human subject's experience on the arrangement that grounds it), it is applied to arrangements that ground subjects, which are maximal by Definition~24.1, so the added condition is met. In the discussion of duplicated brains in Section~24.4 it is applied to two whole brains in separate bodies, each of which is maximal. Theorem~24.6 stands as printed. The general treatment of maximality is part of the boundary problem, which Volume~IV takes up and which MND-01-d treats in a preliminary way.

\section{A terminological slip}

\paragraph{4. Chapter~27, p.~1019, second paragraph.} The sentence \enquote{the phenomenal natures of electrons enter into the grounds of our experience only through their arrangement} should read \enquote{the \emph{protophenomenal} natures of electrons}. The volume holds throughout that the intrinsic natures of fundamental entities are protophenomenal: not experiences, and not such as to make their bearers subjects (Definition~21.2). This is also the position required by MND-18, which shows that a Russellian monism on which those natures are experiences faces a combination problem without solution (MND-18-a, Corollary~3.3). The slip is confined to this sentence; nothing in the surrounding argument uses the stronger reading.

\section{Citation errors}

\paragraph{5. Chapter~21, p.~749, note~3.} The observation that physics describes the world only in terms of structure, and that our acquaintance with the intrinsic character of anything physical is confined to our own experience, is attributed to \enquote{Russell (1921)}. The work cited, \emph{The Analysis of Mind} \citep{Russell1921}, anticipates the point, but the argument referred to is that of \emph{The Analysis of Matter} \citep{Russell1927}; the related discussion of physical knowledge as knowledge of pointer readings is due to \citet{Eddington1928}. The note should cite \citet{Russell1927} and \citet{Eddington1928}, and the entry for \citet{Russell1921} should be retained only for the citation on p.~752, where it is correct.

\paragraph{6. Chapter~22, p.~795, and bibliography.} \citet{Stoljar2001} is listed as appearing in \emph{Mind}. It appeared in \emph{Philosophy and Phenomenological Research}, 62(2). The distinction it draws, between a theory-based and an object-based conception of the physical, is reported correctly in the text.

\section{Effect on the work}

None of the corrections changes the statement of a theorem or corollary. Items~1, 2, 5 and 6 affect references only. Item~4 restores the volume's official terminology in one sentence. Item~3 adds to a lemma a condition that every application of it already satisfies. The thesis of Volume~III, that Russellian monism with protophenomenal natures is a form of materialism in the sense of MND-02 and gives the metaphysics required by Volume~I, stands, as does the conclusion of the work: consciousness is the intrinsic nature of physical structure, and the explanatory gap arises from two modes of presentation of the same facts. Nothing in Volume~I, including its formalization, depends on the corrected passages.

\begin{thebibliography}{9}\small
\bibitem[Eddington(1928)]{Eddington1928} Eddington, A.~S. (1928). \emph{The Nature of the Physical World}. Cambridge University Press.
\bibitem[Russell(1921)]{Russell1921} Russell, B. (1921). \emph{The Analysis of Mind}. George Allen \& Unwin.
\bibitem[Russell(1927)]{Russell1927} Russell, B. (1927). \emph{The Analysis of Matter}. Kegan Paul.
\bibitem[Stoljar(2001)]{Stoljar2001} Stoljar, D. (2001). Two conceptions of the physical. \emph{Philosophy and Phenomenological Research}, 62(2), 253--281.
\end{thebibliography}

\end{document}
